\documentclass[11pt,a4paper]{article}
\usepackage[utf8]{inputenc}
\usepackage[T1]{fontenc}
\usepackage{lmodern}
\usepackage[french]{babel}
\usepackage{amsmath,amssymb,mathtools}
\usepackage{icomma}
\usepackage{xcolor}
\usepackage{tikz}
\usepackage{pgfplots}
\usepackage[most]{tcolorbox}
\usepackage{enumitem}
\usepackage[a4paper,margin=2.2cm,headheight=26pt,headsep=10pt]{geometry}
\usepackage{fancyhdr}
\usepackage[hidelinks]{hyperref}
\usetikzlibrary{arrows.meta,calc}
\pgfplotsset{compat=1.18}
\definecolor{primary}{RGB}{222,49,99}
\definecolor{primarydark}{RGB}{150,22,55}
\definecolor{methcol}{RGB}{52,92,148}
\definecolor{excol}{RGB}{0,135,90}
\definecolor{defcol}{RGB}{0,114,178}
\definecolor{attncol}{RGB}{200,40,55}
\definecolor{curveA}{RGB}{0,90,200}
\definecolor{curveB}{RGB}{210,30,40}
\tcbset{boxbase/.style={enhanced,breakable,boxrule=0.9pt,arc=2.5pt,
  left=3mm,right=3mm,top=2.3mm,bottom=2.3mm,fonttitle=\bfseries,coltitle=white}}
\newtcolorbox{cle}[1][]{boxbase,colback=defcol!6,colframe=defcol,
  title={\textsc{À retenir}\ifx\relax#1\relax\relax\else\textmd{ : \itshape #1}\fi}}
\newtcolorbox{methode}[1][]{boxbase,colback=methcol!7,colframe=methcol,sharp corners,
  title={\textsc{Méthode}\ifx\relax#1\relax\relax\else\textmd{ --- \itshape #1}\fi}}
\newtcolorbox{exemplebox}[1][]{boxbase,colback=excol!5,colframe=excol,
  title={\textsc{Exemple}\ifx\relax#1\relax\relax\else\textmd{ : \itshape #1}\fi}}
\newtcolorbox{piege}[1][]{boxbase,colback=attncol!6,colframe=attncol,
  title={\textsc{Attention}\ifx\relax#1\relax\relax\else\textmd{ : \itshape #1}\fi}}
\newcommand{\R}{\mathbb{R}}
\newcommand{\Z}{\mathbb{Z}}
\newcommand{\ds}{\displaystyle}
\pagestyle{fancy}\fancyhf{}
\fancyhead[L]{\small\textcolor{primarydark}{\textbf{Thème 2} — Logarithmes}}
\fancyhead[R]{\small\textcolor{primarydark}{\textsc{Tutoriel}}}
\fancyfoot[C]{\small\thepage}
\renewcommand{\headrulewidth}{0.8pt}
\begin{document}

\begin{center}
{\LARGE\bfseries\color{primarydark} Thème 2 — Logarithmes}\\[2pt]
{\color{primary}\rule{0.5\textwidth}{1.2pt}}\\[4pt]
{\small\itshape Tutoriel — la mécanique à maîtriser avant les exercices}
\end{center}

\vspace{2mm}
\noindent Le logarithme est la \textbf{fonction réciproque} de l'exponentielle : il retrouve
\emph{l'exposant} à partir du nombre. C'est l'outil pour résoudre toute équation où l'inconnue est
« en exposant ».

\begin{cle}[définition : un logarithme est un exposant]
Pour une base $a>0$, $a\neq1$, et un argument $x>0$ :
\[ \boxed{\ \log_a x = y \iff x=a^{y}\ } \]
$\log_a x$ est « l'exposant qu'il faut donner à $a$ pour obtenir $x$ ». On \emph{ne prend jamais} le
logarithme d'un nombre négatif ou nul. Deux bases ont un nom propre :
$\log x=\log_{10}x$ (décimal) et $\ln x=\log_e x$ (népérien, $e\approx2{,}718$).
\end{cle}

\begin{cle}[les trois propriétés fondamentales]
Pour $x,y>0$, $a>0$ ($a\neq1$), $n\in\Z$ :
\[ \log_a(xy)=\log_a x+\log_a y\ ;\quad
   \log_a\!\Bigl(\frac{x}{y}\Bigr)=\log_a x-\log_a y\ ;\quad
   \log_a\!\bigl(x^{n}\bigr)=n\,\log_a x. \]
\textbf{Idée :} le log \emph{abaisse} d'un cran : produit $\to$ somme, quotient $\to$ différence,
puissance $\to$ produit. C'est ce qui fait « descendre » une inconnue en exposant.
\end{cle}

\begin{cle}[valeurs remarquables et changement de base]
\[ \log_a 1=0\ ;\qquad \log_a a=1\ ;\qquad
   \log_a x=\frac{\log_b x}{\log_b a}\ ;\qquad \log_{a^{n}} x=\frac1n\,\log_a x. \]
En particulier $\ln 1=0$, $\ln e=1$, $\log 10=1$, et $\ln 10\approx 2{,}30$.
\end{cle}

\begin{piege}[ne pas confondre $\ln^{2}x$, $\ln x^{2}$ et $\ln(\ln x)$]
\[ \ln^{2}x=(\ln x)^{2}\ \text{(le log, au carré)}\qquad\neq\qquad \ln x^{2}=2\ln x\ \text{(propriété du log d'un carré)}. \]
Erreur classique aux concours : écrire $\ln^{2}x=2\ln x$. Faux ! Seul $\ln(x^{2})=2\ln x$.
\end{piege}

\begin{minipage}{0.54\textwidth}
\begin{exemplebox}[calculs directs]
\begin{align*}
\log_a\sqrt[n]{x}&=\tfrac1n\log_a x; & \log_a\tfrac1x&=-\log_a x;\\
\ln e^{2}&=2\ln e=2; & \log_{10}10^{2a}&=2a;\\
\ln 0{,}1&=\ln\tfrac{1}{10}=-\ln 10. &&
\end{align*}
\end{exemplebox}
\end{minipage}\hfill
\begin{minipage}{0.42\textwidth}
\centering
\begin{tikzpicture}
\begin{axis}[width=6.1cm,height=4.6cm,axis lines=middle,
 xmin=-2.4,xmax=3.6,ymin=-2.4,ymax=3.6,xtick=\empty,ytick=\empty,
 xlabel={\scriptsize$x$},ylabel={\scriptsize$y$}]
 \addplot[curveA,line width=1.1pt,domain=-2.2:1.2,samples=50]{exp(x)};
 \addplot[curveB,line width=1.1pt,domain=0.06:3.4,samples=60]{ln(x)};
 \addplot[black!40,dashed,domain=-2.2:3.4]{x};
 \node[curveA,font=\scriptsize] at (axis cs:1.2,3.0){$e^{x}$};
 \node[curveB,font=\scriptsize] at (axis cs:3.0,1.35){$\ln x$};
\end{axis}
\end{tikzpicture}\\[-1mm]
{\scriptsize $\ln$ est le symétrique de $\exp$ par $y=x$.}
\end{minipage}

\begin{cle}[monotonie]
Pour $a>1$, $\log_a$ est \emph{croissant} : $\log_a x\geq\log_a y\iff x\geq y$. Pour $0<a<1$, il est
\emph{décroissant} : l'inégalité \emph{change de sens}. C'est essentiel pour les inéquations.
\end{cle}

\begin{methode}[résoudre une équation « en exposant » ou logarithmique]
\begin{enumerate}[leftmargin=1.8em,topsep=1pt,itemsep=1pt]
  \item \textbf{Domaine d'abord (C.E.)} : tout argument de log doit être $>0$.
  \item inconnue \emph{en exposant} : appliquer $\ln$ aux deux membres, puis $\ln(a^{u})=u\ln a$ ;
  \item inconnue \emph{dans un log} : repasser à l'exponentiel, $\log_a u=v\iff u=a^{v}$ ;
  \item résoudre, puis \textbf{vérifier} que les solutions respectent la C.E.
\end{enumerate}
\end{methode}

\begin{exemplebox}[deux résolutions types]
\textbf{(1)} $e^{2x}=10\Rightarrow \ln e^{2x}=\ln 10\Rightarrow 2x=\ln 10\Rightarrow x=\tfrac{\ln 10}{2}.$\\
\textbf{(2)} $\ln(x-1)\geq\ln 2$. C.E. : $x>1$. $\ln$ croissant $\Rightarrow x-1\geq 2\Rightarrow x\geq 3$.
Avec la C.E. : $S=[3,+\infty[$.
\end{exemplebox}

\end{document}
